\newcounter{english}
\newcounter{stoppagebreaks}
\documentclass{article}

% packages
\usepackage{amsmath, amsthm, thmtools, amsfonts, amssymb, luacode, catchfile, hyperref, ifthen, graphicx, bbm}
\ifcsname c@kobocompile\endcsname
	\usepackage[a5paper, total={1072pt, 1448pt}, margin=10pt, includeheadfoot]{geometry} % set page margins
\else
	\usepackage[a4paper, margin=50pt, includeheadfoot]{geometry}
\fi
\usepackage[shortlabels]{enumitem}
\usepackage[skip=3pt, indent=0pt]{parskip}

% language
\usepackage[bidi=basic, layout=tabular, provide=*]{babel}
\ifcsname c@english\endcsname
	\babelprovide[main, import]{english}
\else
	\babelprovide[main, import, alph=letters.gershayim]{hebrew}
	\setlist[enumerate,2]{label={\alph*.}}
	\babelprovide{rl}
\fi
\babelfont{rm}[Renderer=Harfbuzz]{Libertinus Serif}
\babelfont{sf}{Libertinus Sans}
\babelfont{tt}{Libertinus Mono}

% style
\ifcsname c@stoppagebreaks\endcsname
\else
	\AddToHook{cmd/section/before}{\clearpage}	% Add line break before section
\fi
\linespread{1.3}
\setcounter{secnumdepth}{0}		% Remove default number tags from sections, this won't do well with theorems
\AtBeginDocument{\setlength{\belowdisplayskip}{3pt}}
\AtBeginDocument{\setlength{\abovedisplayskip}{3pt}}
\graphicspath{ {../images/} }

% operators
\DeclareMathOperator\cis{cis}
\DeclareMathOperator\Sp{Span}
\DeclareMathOperator\tr{tr}
\DeclareMathOperator\im{Im}
\DeclareMathOperator\re{Re}
\DeclareMathOperator\diag{diag}
\DeclareMathOperator*\lowlim{\underline{lim}}
\DeclareMathOperator*\uplim{\overline{lim}}
\DeclareMathOperator\rng{rng}
\DeclareMathOperator\Sym{Sym}
\DeclareMathOperator\Arg{Arg}
\DeclareMathOperator\Log{Log}
\DeclareMathOperator\dom{dom}
\DeclareMathOperator\supp{Supp}
\DeclareMathOperator\var{Var}
\DeclareMathOperator\cov{Cov}
\DeclareMathOperator\End{End}
\DeclareMathOperator\id{id}
\DeclareMathOperator\cl{cl}
\DeclareMathOperator\gal{Gal}
\DeclareMathOperator\dist{dist}
\DeclareMathOperator\form{form}
\DeclareMathOperator\term{term}
\DeclareMathOperator\sent{sent}
\DeclareMathOperator\Th{Th}
\DeclareMathOperator\Hom{Hom}
\DeclareMathOperator\Aut{Aut}
\DeclareMathOperator\Char{char}
\DeclareMathOperator\vol{vol}
\DeclareMathOperator\GL{GL}
\DeclareMathOperator\SL{SL}
\DeclareMathOperator\lcm{lcm}

% commands
\newcommand{\cupdot}{\mathbin{\mathaccent\cdot\cup}}
\newcommand{\indicator}{\mathbbm{1}}
\newcommand{\AAA}[0]{\mathbb{A}}
\newcommand{\Aa}[0]{\mathcal{A}}
\newcommand{\BB}[0]{\mathbb{B}}
\newcommand{\Bb}[0]{\mathcal{B}}
\newcommand{\CC}[0]{\mathbb{C}}
\newcommand{\Cc}[0]{\mathcal{C}}
\newcommand{\DD}[0]{\mathbb{D}}
\newcommand{\Dd}[0]{\mathcal{D}}
\newcommand{\EE}[0]{\mathbb{E}}
\newcommand{\FF}[0]{\mathbb{F}}
\newcommand{\Ff}[0]{\mathcal{F}}
\newcommand{\Ii}[0]{\mathcal{I}}
\newcommand{\Gg}[0]{\mathcal{G}}
\newcommand{\HH}[0]{\mathbb{H}}
\newcommand{\Hh}[0]{\mathcal{H}}
\newcommand{\Kk}[0]{\mathcal{K}}
\newcommand{\Ll}[0]{\mathcal{L}}
\newcommand{\Mm}[0]{\mathcal{M}}
\newcommand{\NN}[0]{\mathbb{N}}
\newcommand{\Nn}[0]{\mathcal{N}}
\newcommand{\PP}[0]{\mathbb{P}}
\newcommand{\Pp}[0]{\mathcal{P}}
\newcommand{\QQ}[0]{\mathbb{Q}}
\newcommand{\RR}[0]{\mathbb{R}}
\newcommand{\Rr}[0]{\mathcal{R}}
\newcommand{\Ss}[0]{\mathcal{S}}
\newcommand{\TT}[0]{\mathbb{T}}
\newcommand{\Tt}[0]{\mathcal{T}}
\newcommand{\UU}[0]{\mathbb{U}}
\newcommand{\Uu}[0]{\mathcal{U}}
\newcommand{\Vv}[0]{\mathcal{V}}
\newcommand{\Ww}[0]{\mathcal{W}}
\newcommand{\Xx}[0]{\mathcal{X}}
\newcommand{\Yy}[0]{\mathcal{Y}}
\newcommand{\ZZ}[0]{\mathbb{Z}}
\newcommand{\Zz}[0]{\mathcal{Z}}
\newcommand{\acts}[0]{\circlearrowright}
\newcommand{\explain}[2] {
	\begin{flalign*}
		 && \text{#2} && \text{#1}
	\end{flalign*}
}
\newcommand{\maketitleprint}[1][green]{
	\begin{center}
		\includegraphics[width=6cm]{dragon_#1}
	\end{center}
}

% theorem commands
\newtheoremstyle{c_remark}
	{}	% Space above
	{}	% Space below
	{}% Body font
	{}	% Indent amount
	{\bfseries}	% Theorem head font
	{}	% Punctuation after theorem head
	{.5em}	% Space after theorem head
	{\thmname{#1}\thmnumber{ #2}\thmnote{ \normalfont{\text{(#3)}}}}	% head content
\newtheoremstyle{c_definition}
	{3pt}	% Space above
	{3pt}	% Space below
	{}% Body font
	{}	% Indent amount
	{\bfseries}	% Theorem head font
	{}	% Punctuation after theorem head
	{.5em}	% Space after theorem head
	{\thmname{#1}\thmnumber{ #2}\thmnote{ \normalfont{\text{(#3)}}}}	% head content
\newtheoremstyle{c_plain}
	{3pt}	% Space above
	{3pt}	% Space below
	{\itshape}% Body font
	{}	% Indent amount
	{\bfseries}	% Theorem head font
	{}	% Punctuation after theorem head
	{.5em}	% Space after theorem head
	{\thmname{#1}\thmnumber{ #2}\thmnote{ \text{(#3)}}}	% head content

\ifcsname c@english\endcsname
	\theoremstyle{plain}
	\newtheorem{theorem}{Theorem}[section]
	\newtheorem*{theorem*}{Theorem}
	\newtheorem{lemma}[theorem]{Lemma}
	\newtheorem{proposition}[theorem]{Proposition}
	\newtheorem*{proposition*}{Proposition}

	\theoremstyle{definition}
	\newtheorem{definition}[theorem]{Definition}
	\newtheorem*{definition*}{Definition}
	\newtheorem{example}{Example}[section]
	\newtheorem{exercise}{Exercise}[section]

	\theoremstyle{remark}
	\newtheorem*{remark}{Remark}
	\newtheorem*{solution}{Solution}
	\newtheorem{corollary}[theorem]{Corollary}
	\newtheorem{notation}[theorem]{Notation}
\else
	\theoremstyle{c_plain}
	\newtheorem{theorem}{משפט}[section]
	\newtheorem*{theorem*}{משפט}
	\newtheorem{lemma}[theorem]{למה}
	\newtheorem{proposition}[theorem]{טענה}
	\newtheorem*{proposition*}{טענה}

	\theoremstyle{c_definition}
	\newtheorem{definition}[theorem]{הגדרה}
	\newtheorem*{definition*}{הגדרה}
	\newtheorem{example}{דוגמה}[section]
	\newtheorem{exercise}{תרגיל}[section]

	\theoremstyle{c_remark}
	\newtheorem*{remark}{הערה}
	\newtheorem*{solution}{פתרון}
	\newtheorem{corollary}[theorem]{מסקנה}
	\newtheorem{notation}[theorem]{סימון}
\fi

% Questions related commands
\newcounter{question}
\setcounter{question}{1}
\newcounter{sub_question}
\setcounter{sub_question}{1}

\ifcsname c@english\endcsname
	\newcommand{\question}[1][0]{
		\ifthenelse{#1 = 0}{}{\setcounter{question}{#1}}
		\section{Question \arabic{question}}
		\addtocounter{question}{1}
		\setcounter{sub_question}{1}
	}

	\newcommand{\subquestion}[1][0]{
		\ifthenelse{#1 = 0}{}{\setcounter{sub_question}{#1}}
		\subsection{Part \alph{sub_question}}
		\addtocounter{sub_question}{1}
	}
\else
	\newcommand{\question}[1][0]{
		\ifthenelse{#1 = 0}{}{\setcounter{question}{#1}}
		\section{שאלה \arabic{question}}
		\addtocounter{question}{1}
		\setcounter{sub_question}{1}
	}

	\newcommand{\subquestion}[1][0]{
		\ifthenelse{#1 = 0}{}{\setcounter{sub_question}{#1}}
		\subsection{סעיף \localecounter{letters.gershayim}{sub_question}}
		\addtocounter{sub_question}{1}
	}
\fi

% import lua and start of document
\directlua{common = require ('../common')}

\GetEnv{AUTHOR}

% headers
\author{\AUTHOR}
\date\today

\title{Lie Groups --- Lecture 3}
\setcounter{secnumdepth}{2}
% chktex-file 9
% chktex-file 17

\usepackage{fancyhdr}
\pagestyle{fancy}
\renewcommand{\headrulewidth}{0pt}

\begin{document}
\maketitle

\section{Reminder}
\begin{definition}[Tangent space]
	If $(X, A)$ is a smooth manifold and $p \in X$ is some point, then the tangent space of $p$ is defined as,
	\[
		T_p X
		= \{ \gamma : (-\delta, \delta) \to X \mid \gamma(0) = p, \gamma \in C^{\infty}(X), \delta \in \RR^+ \} / \sim
	\]
\end{definition}
\begin{definition}[exponential map]
	If $G$ is a Lie group then define $\exp : L(G) \to G$ as the map that for any $v \in L(G)$, $\theta_v$ be the corresponding $1$-parameter subgroup of $G$, then $\exp tv = \theta_v(t)$.
\end{definition}

\section{Submanifolds}
\begin{definition}[Submanifold]
	Let $M$ be a smooth manifold modeled on $V$, and let $N \subseteq M$ be a subset and $W \subseteq V$ be a subspace.
	$N$ is a \textbf{submanifold} modeled on $W$ if for any $x \in N$ there exists a chart compatible with $M$, $\varphi_{\alpha} : M_{\alpha} \to \Omega_{\alpha}$
	where $M_{\alpha} \subseteq M, \Omega_{\alpha} \subseteq V$ are open, that restricts to a bijection $M_{\alpha} \cap N \to \Omega_{\alpha} \cap W$.
\end{definition}
\begin{proposition}
	Given a manifold $M$ and a submanifold $N$, the bijections $\{ M_{\alpha} \cap N \to \Omega_{\alpha} \cap W \}$ are an atlas on $N$. \\
	Furthermore, equivalent atlases on $M$ yield equivalent atlases on $N$ thus $N$ is a smooth manifold modeled on $W$.
\end{proposition}
\begin{proof}
	Left as an exercise.
\end{proof}
\begin{proposition}
	If $N$ is a submanifold of $M$, then $N$ can be given differential structure as a manifold so that:
	\begin{enumerate}
		\item The inclusion of $N$ in $M$ is smooth
		\item $P \to N$ is smooth if and only if the composition $P \to N \to M$ is smooth
	\end{enumerate}
\end{proposition}
\begin{proof}
	Let $f : N \to M$ be the identity embedding, and let $\varphi : U \to V$ be a chart of $M$, $u \subseteq M, V \subseteq \RR^k$ are open, and assume $N \cap U \ne \emptyset$.
	Then $\varphi$ can be restricted to a chart $\psi : N \cap U \to V \cap W$ which is a bijection and $\varphi \circ f \circ \psi$ is smooth by our definition of smoothness.

	The second part is similar and is left as an exercise.
\end{proof}
\begin{corollary}
	$T(N)$ is embedded in $T(M)$.
\end{corollary}
\begin{exercise}
	Show that if $N_1, N_2$ are manifolds of $M_1, M_2$ then $N_1 \times N_2$ is a submanifold of $M_1 \times M_2$ and its differential structure as a submanifold is the same as of the product.
\end{exercise}
\begin{proposition}
	If $G$ is a Lie group and $H$ is both a submanifold and a subgroup then $H$ is a Lie group.
\end{proposition}
\begin{proof}
	Follows from the last exercise and proposition for the maps $\mu : G \times G \to G$ and $\mu : H \times H \to H$ and $i : G \to G$ to $i : H \to H$.
\end{proof}

\section{Closed subgroups}
\begin{definition}[Closed subgroup]
	A subgroup $H$ of a Lie group $G$ is closed if it is closed as a subset of the topology $G$.
\end{definition}
\begin{theorem}[Closed subgroup as submanifold]
	A closed subgroup $H$ of a Lie group is a submanifold.
\end{theorem}
To prove the theorem is given by the next three lemmas.
\begin{lemma}
	Suppose $L(G) = T_e G$ has a norm and suppose that $0 \ne h_n \in L(G)$ is a sequence of points such that,
	\[
		\forall n \in \NN,\ \exp h_n \in H,
		\quad
		h_n \xrightarrow{n \to \infty} 0,
		\quad
		\frac{1}{|h_n|} h_n \to v \in L(G)
	\]
	Then $\exp(t v) \in H$ for all $t \in \RR$.
\end{lemma}
In other words, a closed subgroup is closed under ``directions''.
\begin{proof}
	TODO
\end{proof}
\begin{lemma}
	Let $W$ be a vector subspace of $L(G)$, then $W$ is a vector subspace of $L(G)$.
\end{lemma}
\begin{proof}
	TODO
\end{proof}
\begin{lemma}
	$\exp W$ is a neighborhood of $e \in H$.
\end{lemma}
\begin{proof}
	TODO
\end{proof}
\begin{proof}[Proof of the theorem]
	By the last lemma $\exp W$ is a neighborhood of $e$ in $H$, thus $\exp |_W$ is a chart of a neighborhood of $e$.
	This chart can be translated to any other point $\in H$ by multiplication.
\end{proof}
\begin{example}
	$O(n) \subseteq \GL_n(\RR)$ the subspace of orthogonal invertible matrices is a closed subgroup (TODO proof) and so a submanifold of a Lie group thus its is a Lie group.
\end{example}
\begin{example}
	$U(n) \subseteq \GL_n(\CC)$ the subspace of unitary complex matrices is also a closed subset (TODO proof) thus also a Lie group.
\end{example}
TODO\@: Should we copy the third example?
\begin{example}
	If $G$ is a compact Lie group and $H$ is a closed connected Abelian subgroup, then $H$ is a torus.
\end{example}
\begin{proposition}
	A closed connected subgroup $H$ of a Lie group $G$ is determined by $L(G)$.
\end{proposition}
\begin{proof}
	TODO something with 2.17 which means that a homomorphism is determined by its projection to the tangent space.
\end{proof}

\section{Quotient spaces}
\begin{definition}[Subgroup \& Cosets]
	Suppose $G$ is a Lie group and $H$ is a closed subgroup, then the \textbf{quotient space} $G / H$ is the set of (left) cosets $\{ g H \mid g \in G \}$.
	The map $\pi : G \to G / H$ given by $g \mapsto g H$ is the projection, and give $G / H$ the quotient topology.
\end{definition}
\begin{proposition}
	$G / H$ is Hausdorff.
\end{proposition}
\begin{proof}
	Suppose that $x H, y H \in G / H$ are different points, conversely $x, y \in G$ such that $x \ne y$ and $x H \ne y H$.
	$G$ is Hausdorff then there are $x \in U, y \in V$ open neighborhoods such that $U \cap V = \emptyset$.
	$e \in x^{-1} U, y^{-1} V$ then there is a chart $\varphi : W \to \RR^k$ for $W \subseteq x^{-1} U \cap y^{-1} V$ up to compatibility with $G$.
	By translating the chart we can make $V, U$ arbitrarily small (taking $\varphi^{-1}(B(0, \varepsilon))$), such that $\varepsilon < \dist(\varphi(x W \cap x H, y W \cap y H))$ if it is defined (otherwise the claim holds).
\end{proof}
Together with the fact that a quotient topology assumes 2nd countability of a topology, we get,
\begin{proposition}
	$G / H$ can be given a smooth manifold structure, such that $\pi$ is smooth and any $f : G / H \to M$ is smooth if and only if there exists $F : G \to M$ smooth with $f = F \circ \pi$.
\end{proposition}
\begin{proof}
	TODO
\end{proof}
\begin{definition}[Bundle]
	TODO
\end{definition}
\begin{proposition}
	$H \to G \to G / H$ is a bundle.
\end{proposition}

\end{document}
